\begin{description}
\item[(a)]
  \[
    \theta_{\mathrm{E}}
    = \sqrt{\frac{4GM}{c^2}\,\frac{D_s - D_\ell}{D_s D_\ell}}
    = \sqrt{\frac{4GM}{\mathrm{au}\,c^2}
      \left(\frac{\mathrm{au}}{D_\ell} - \frac{\mathrm{au}}{D_s}\right)}
    = 2.7\,\mathrm{mas}
  \]
  The angular resolution of modern large ($D \approx 10$\,m) optical
  ($\lambda = 550$\,nm) telescopes is
  $\theta_0 = 1.22\lambda/D \approx 14$\,mas. Thus,
  $\theta_0 \gg \theta_{\mathrm{E}}$, so the images created during
  microlensing events cannot be resolved by these telescopes.
  \hfill(2 points)

\item[(b)]
  \begin{align*}
    \theta_c &= \frac{\theta_+ A_+ + \theta_- A_-}{A_+ + A_-}
    = \frac{\frac{1}{4}\left(u + \sqrt{u^2+4}\right)
        \left(\frac{u^2+2}{u\sqrt{u^2+4}} + 1\right)
      + \frac{1}{4}\left(u - \sqrt{u^2+4}\right)
        \left(\frac{u^2+2}{u\sqrt{u^2+4}} - 1\right)}
      {\frac{u^2+2}{u\sqrt{u^2+4}}}\,\theta_{\mathrm{E}}\\
    &= \frac{\left(u + \sqrt{u^2+4}\right)
        \left(u^2 + 2 + u\sqrt{u^2+4}\right)
      + \left(u - \sqrt{u^2+4}\right)
        \left(u^2 + 2 - u\sqrt{u^2+4}\right)}
      {4(u^2+2)}\,\theta_{\mathrm{E}}\\
    &= \frac{2u(u^2+2) + 2u(u^2+4)}{4(u^2+2)}\,\theta_{\mathrm{E}}
    = \frac{2u(2u^2+6)}{4(u^2+2)}\,\theta_{\mathrm{E}}
    = \frac{u(u^2+3)}{u^2+2}\,\theta_{\mathrm{E}}
  \end{align*}
  \hfill(8 points)

\item[(c)]
  \[
    \Delta\theta = \theta_c - \theta
    = \frac{u(u^2+3)}{u^2+2}\,\theta_{\mathrm{E}} - u\theta_{\mathrm{E}}
    = \frac{u(u^2+3) - u(u^2+2)}{u^2+2}\,\theta_{\mathrm{E}}
    = \frac{u}{u^2+2}\,\theta_{\mathrm{E}}
  \]
  \hfill(3 points)
  \[
    \Delta\theta(u = 0) = 0
  \]
  so there is no deflection when the lens and the source are nearly
  perfectly aligned. \hfill(1 points)

\item[(d)] Plots of the measured E and N positions against time.
  \hfill(10 points, 5 points for each graph)
  % FIGURE NEEDED: two scatter plots with error bars, Position E (mas) and
  % Position N (mas) vs Time (HJD-2450000), from 2023 DA solutions page 48

\item[(e)] I will use the fact that the first epoch of astrometric
  observations was taken close to the peak of the light curve (that is,
  $u_1 \approx 0$, that is, almost no astrometric deflection). Similarly,
  astrometric deflection is close to zero for the last epoch. Thus,
  \begin{align*}
    \mu_E &\approx \frac{x_{E,8} - x_{E,1}}{t_8 - t_1}
      = -2.247 \pm 0.029\,\mathrm{mas/yr}\\
    \mu_N &\approx \frac{x_{N,8} - x_{N,1}}{t_8 - t_1}
      = -3.674 \pm 0.038\,\mathrm{mas/yr}
  \end{align*}
  \hfill(8 points)

  \emph{No points should be given if a student does not recognize that the
  deflection is zero during the first epoch, e.g., they try fitting a
  straight line to all data points.}

\item[(f)] I fitted a straight line joining the first and last epoch data
  and then subtracted it from astrometric measurements. This is because the
  observed path of the source on the sky is the sum of two effects: the
  rectilinear proper motion of the source and the astrometric deflection:
  \begin{align*}
    x(E) &= x_1^E + (t - t_1)\mu_{\mathrm{E}} + \Delta\theta(E)\\
    x(N) &= x_1^N + (t - t_1)\mu_{\mathrm{N}} + \Delta\theta(N),
  \end{align*}
  where $x_1^E$ is the East position of the source during the first epoch,
  $x_1^N$ is the North position of the source during the first epoch, $t_1$
  is the time of the first observation, and $\Delta\theta(E)$ and
  $\Delta\theta(N)$ is the astrometric deflection due to microlensing in
  East and North direction, respectively.
  % FIGURE NEEDED: Figure 7 -- E and N position vs time with blue lines
  % joining first and last epoch data points (2023 DA solutions page 49)

  Thus, the astrometric deflection during $i$th epoch is:
  \begin{align*}
    \Delta\theta(E)_i &= x_i^E - x_1^E - (t_i - t_1)\mu_{\mathrm{E}}\\
    \Delta\theta(N)_i &= x_i^N - x_1^N - (t_i - t_1)\mu_{\mathrm{N}}
  \end{align*}
  and the total deflection is:
  \[
    \Delta\theta_i = \sqrt{\Delta\theta(E)_i^2 + \Delta\theta(N)_i^2}.
  \]
  I will also use the fact that $u_0 \approx 0$, so
  $u = (t - t_0)/t_E$, where $t_0 = 2455763$ is the peak time. Results of
  my calculations are shown in the table below.
  % FIGURE NEEDED: Figure 8 -- deflection E and N vs time (not required of
  % students); Figure 9 -- total astrometric deflection vs u, the plot
  % students are asked to make (2023 DA solutions page 50)

  \begin{center}
  \begin{tabular}{rrrrr}
  \hline\hline
  Epoch & $u$ & $\Delta\theta$ (E,mas) & $\Delta\theta$ (N,mas)
    & $\Delta\theta$ (mas)\\
  \hline
  1 & 0.01 & $0.00 \pm 0.13$ & $-0.00 \pm 0.16$ & $0.00 \pm 0.21$\\
  2 & 0.42 & $0.36 \pm 0.12$ & $-0.84 \pm 0.24$ & $0.91 \pm 0.27$\\
  3 & 1.69 & $0.43 \pm 0.14$ & $-1.50 \pm 0.08$ & $1.56 \pm 0.16$\\
  4 & 1.75 & $0.95 \pm 0.36$ & $-1.40 \pm 0.77$ & $1.69 \pm 0.85$\\
  5 & 2.69 & $0.47 \pm 0.21$ & $-1.58 \pm 0.12$ & $1.65 \pm 0.24$\\
  6 & 3.34 & $0.44 \pm 0.07$ & $-0.90 \pm 0.40$ & $1.00 \pm 0.41$\\
  7 & 4.83 & $0.21 \pm 0.25$ & $-0.47 \pm 0.29$ & $0.51 \pm 0.38$\\
  8 & 9.04 & $0.00 \pm 0.12$ & $-0.00 \pm 0.17$ & $0.00 \pm 0.21$\\
  \hline
  \end{tabular}
  \end{center}
  \hfill(20 points)

\item[(g)] We would like to fit the function
  $\Delta\theta = \frac{u}{u^2+2}\theta_{\mathrm{E}}$ to the data. Thus, my
  ``new'' independent variable would be $x' = u/(u^2+2)$. Now, I would like
  to fit the function $y' = \theta_{\mathrm{E}} x'$, where
  $y' = \Delta\theta$. Thus
  \[
    \theta_{\mathrm{E}}
    = \frac{\sum y'_i x'_i/\sigma_i^2}{\sum x_i^{\prime 2}/\sigma_i^2}
      \pm \frac{1}{\sqrt{\sum x_i^{\prime 2}/\sigma_i^2}}
    = 4.5 \pm 0.4\,\mathrm{mas}
  \]
  % FIGURE NEEDED: deflection (mas) vs u/(u^2+2) with fitted straight line
  % (2023 DA solutions page 51)

  Alternative solution: The gravitational deflection is described by the
  formula $\Delta\theta = \frac{u}{u^2+2}\theta_{\mathrm{E}}$. It can be
  demonstrated that this function reaches a local maximum for
  $u = \sqrt{2}$ with
  $\Delta\theta_{\mathrm{max}} = \frac{\sqrt{2}}{4}\theta_{\mathrm{E}}
  = 0.354\,\theta_{\mathrm{E}}$. From the data, we can estimate the maximum
  deflection of $\Delta\theta_{\mathrm{max}} = 1.59 \pm 0.14$\,mas. Thus,
  $\theta_{\mathrm{E}} = 4.5 \pm 0.4$\,mas. \hfill(16 points)

\item[(h)] Let $\pi_{\mathrm{rel}} = \pi_l - \pi_s$ be the relative
  lens--source parallax. From the definition of the angular Einstein radius
  we have
  \[
    \theta_{\mathrm{E}}
    = \sqrt{\frac{4GM}{c^2}\,\frac{D_s - D_l}{D_s D_l}}
    = \sqrt{\frac{4GM}{c^2}\left(\frac{1}{D_l} - \frac{1}{D_s}\right)}
    = \sqrt{\frac{4GM}{c^2\,\mathrm{au}}
      \left(\frac{\mathrm{au}}{D_l} - \frac{\mathrm{au}}{D_s}\right)}
    = \sqrt{\frac{4GM\pi_{\mathrm{rel}}}{c^2\,\mathrm{au}}}.
  \]
  From the definition of the microlensing parallax
  $\pi_{\mathrm{rel}} = \pi_{\mathrm{E}}\theta_E$, so this equation
  becomes:
  \[
    \theta_{\mathrm{E}}
    = \sqrt{\frac{4GM\pi_{\mathrm{E}}\theta_{\mathrm{E}}}{c^2\,\mathrm{au}}},
  \]
  and hence
  \[
    M = \frac{\theta_{\mathrm{E}} c^2\,AU}{4G\pi_{\mathrm{E}}}
    = 5.8 \pm 0.8\,M_\odot.
  \]
  The uncerainty% sic
  {} on $M$ can be determined from the relation:
  \[
    \frac{\Delta M}{M}
    = \sqrt{\left(\frac{\Delta\theta_E}{\theta_{\mathrm{E}}}\right)^2
      + \left(\frac{\Delta\pi_{\mathrm{E}}}{\pi_{\mathrm{E}}}\right)^2}
  \]
  \hfill(7 points)
\end{description}
